\section{Implementing Dirichlet data}\label{app:dirichlet}

Strict Dirichlet data at an artificial cut cannot be imposed simultaneously with arbitrary field matching: if both cut traces are fixed to zero from the outset, their difference carries no gluing information. The useful construction is therefore an interpolation. A boundary penalty realizes decoupled Dirichlet theories at one endpoint, while a relative-trace penalty glues the regions at the other endpoint. The scalar interpolation is diagonalized numerically in Section $\ref{subsec:numerics-interval}$.

\subsection{Scalar boundary penalty}\label{subsec:dirichlet-scalar}

For the scalar interval, keep the physical outer Dirichlet data exact and vary the two traces at the artificial cut. The reflection-symmetric local quadratic form contains two independent positive coefficients,

\begin{align}
S_{d,h}^{\mathrm D}
=\sum_{i=1}^2S_i
-\frac12\int\mathrm dt\left[
d(\phi_1^2+\phi_2^2)
+h(\phi_1-\phi_2)^2
\right]_{x=0},
\qquad d,h>0.
\end{align}

The coefficient $d$ is the boundary penalty and $h$ is the gluing interaction. Their endpoint requirements are different:

\begin{align}
\text{Dirichlet endpoint}:\quad d\to\infty,\quad h\to0,
\end{align}

\begin{align}
\text{smooth-gluing endpoint}:\quad d\to0,\quad h\to\infty.
\end{align}

Thus two boundary operators are required. A convenient one-parameter trajectory through their two-dimensional coupling space is obtained by introducing a fixed inverse-length scale $\mu>0$ and setting

\begin{align}
d_\kappa=\frac\mu\kappa,
\qquad
h_\kappa=\mu\kappa,
\qquad \kappa>0.
\end{align}

Along this continuum trajectory the action is

\begin{align}
S_\kappa^{\mathrm D}
=\sum_{i=1}^2 S_i
-\frac12\int\mathrm dt\left[
\frac\mu\kappa(\phi_1^2+\phi_2^2)
+\mu\kappa(\phi_1-\phi_2)^2
\right]_{x=0}.
\end{align}

Free variation of the cut traces gives

\begin{align}
\phi_1'
+\frac\mu\kappa\phi_1
+\mu\kappa(\phi_1-\phi_2)=0,
\end{align}

\begin{align}
\phi_2'
-\frac\mu\kappa\phi_2
+\mu\kappa(\phi_1-\phi_2)=0,
\end{align}

at $x=0$. The positive spatial quadratic form is

\begin{align}
Q_\kappa^{\mathrm D}
=\sum_{i=1}^2\int_{I_i}\mathrm dx\,
(|\phi_i'|^2+m^2|\phi_i|^2)
+\left[
\frac\mu\kappa(|\phi_1|^2+|\phi_2|^2)
+\mu\kappa|\phi_1-\phi_2|^2
\right]_{x=0}.
\end{align}

For a family with bounded $Q_\kappa^{\mathrm D}$,

\begin{align}
\kappa\to0:
\qquad
\phi_1(0)=\phi_2(0)=0.
\end{align}

This is the direct sum of the two strict Dirichlet regions. In the opposite limit,

\begin{align}
\kappa\to\infty:
\qquad
\phi_1(0)=\phi_2(0).
\end{align}

Subtracting the two finite-$\kappa$ cut equations then gives

\begin{align}
\phi_1'(0)=\phi_2'(0),
\end{align}

so the fields join into one scalar on the full interval. Because neither cut term contains a time derivative, the pre-symplectic form remains the sum of the two bulk forms throughout the interpolation.

The two operators are essential. The symbol $\kappa$ only parameterizes the selected continuum trajectory $d_\kappa h_\kappa=\mu^2$; it does not mean that one boundary operator realizes both limits.

This distinction becomes unavoidable after a mode cutoff. In the common and relative sectors the continuum stiffnesses are

\begin{align}
\alpha_+=d,
\qquad
\alpha_-=d+2h.
\end{align}

If $T_N(0)$ is the omitted zero-energy boundary response, low-energy matching requires

\begin{align}
\frac1{\alpha_{\sigma,N}}
=\frac1{\alpha_\sigma}+T_N(0),
\qquad \sigma=\pm,
\end{align}

and therefore

\begin{align}
d_N=\alpha_{+,N},
\qquad
h_N=\frac{\alpha_{-,N}-\alpha_{+,N}}2.
\end{align}

On the continuum trajectory $d h=\mu^2$, these matched coefficients obey

\begin{align}
d_Nh_N
=\frac{\mu^2}{
\left(1+\alpha_-T_N(0)\right)
\left(1+\alpha_+T_N(0)\right)^2},
\end{align}

which differs from $\mu^2$ whenever the omitted response is nonzero. Hence there is generally no single bare $\kappa_N$ satisfying both $d_N=\mu/\kappa_N$ and $h_N=\mu\kappa_N$. The continuum one-parameter path remains a valid family of exact boundary conditions, but it is not closed under finite-cutoff response matching.

\subsection{Gauge-invariant Maxwell boundary penalty}\label{subsec:dirichlet-maxwell}

For Maxwell theory, Dirichlet data fix the tangential connection only up to a gauge transformation. On the two cut faces $\Gamma_i$, introduce boundary Stueckelberg fields $\varphi_i$ with

\begin{align}
A_i\longrightarrow A_i+\mathrm d\lambda_i,
\qquad
\varphi_i\longrightarrow\varphi_i+\lambda_i\big|_{\Gamma_i},
\end{align}

and define the dressed tangential fields

\begin{align}
a_{i,a}=A_{i,a}-\partial_a\varphi_i,
\qquad
a=t,y.
\end{align}

Along the same convenient continuum trajectory, the gauge-invariant Dirichlet penalty is

\begin{align}
S_{\mathrm D}^{(\kappa)}
=-\frac\mu{2\kappa}
\sum_{i=1}^2\int_{\Gamma_i}\mathrm dt\mathrm dy\,
\gamma^{ab}a_{i,a}a_{i,b}.
\end{align}

Together with the Maxwell bulk variation, it gives

\begin{align}
n_{i,\mu}F_i^{\mu a}
+\frac\mu\kappa a_i^a=0,
\qquad
\partial_a a_i^a=0.
\end{align}

The strict limit $\kappa\to0$ forces $a_{i,a}\to0$ for bounded energy, which is precisely the gauge-invariant Dirichlet condition on each cut face.

To interpolate from these two Dirichlet regions to the uncut theory along this trajectory, add

\begin{align}
S_{\mathrm{glue}}^{(\kappa)}
=-\frac{\mu\kappa}{2}\int_\Gamma\mathrm dt\mathrm dy\,
\gamma^{ab}(a_{1,a}-a_{2,a})(a_{1,b}-a_{2,b}).
\end{align}

More generally the two coefficients are independent. With

\begin{align}
B_{d,h}
=\begin{pmatrix}
d+h&-h\\
-h&d+h
\end{pmatrix},
\end{align}

the common and relative eigenvalues are $d$ and $d+2h$. The one-parameter continuum choice is $B_{d_\kappa,h_\kappa}=\mu M_\kappa$, with

\begin{align}
M_\kappa=
\begin{pmatrix}
\kappa^{-1}+\kappa&-\kappa\\
-\kappa&\kappa^{-1}+\kappa
\end{pmatrix},
\end{align}

and the complete cut action along that trajectory is

\begin{align}
S_{\partial,\kappa}
=-\frac\mu2\int_\Gamma\mathrm dt\mathrm dy\,
\gamma^{ab}\boldsymbol a_a^{\mathrm T}M_\kappa\boldsymbol a_b.
\end{align}

The cut equations are

\begin{align}
n_{i,\mu}F_i^{\mu a}
+\mu\sum_{j=1}^2(M_\kappa)_{ij}a_j^a=0,
\end{align}

\begin{align}
\partial_a\left[
\sum_{j=1}^2(M_\kappa)_{ij}a_j^a
\right]=0.
\end{align}

The variation of the boundary fields also has an endpoint term in time. Its contribution to $\theta_\kappa$ is

\begin{align}
\begin{aligned}
\theta_\kappa
&=\sum_{i=1}^2\int_{\Sigma_i}\mathrm dx\mathrm dy\,
(E_{i,x}\delta A_{i,x}+E_{i,y}\delta A_{i,y}) \\
&\quad+\mu\int_{S_y^1}\mathrm dy\,
\sum_{i,j=1}^2(M_\kappa)_{ij}
(\dot\varphi_j-A_{j,t})\delta\varphi_i.
\end{aligned}
\end{align}

Consequently, the boundary fields enter the CPS form even though they do not supply independent physical oscillators after the gauge quotient.

In the basis

\begin{align}
a_{\pm,a}=\frac{a_{1,a}\pm a_{2,a}}{\sqrt2},
\end{align}

the cut action becomes

\begin{align}
S_{\partial,\kappa}
=-\frac\mu2\int_\Gamma\mathrm dt\mathrm dy\,
\gamma^{ab}\left[
\frac1\kappa a_{+,a}a_{+,b}
+\left(\frac1\kappa+2\kappa\right)a_{-,a}a_{-,b}
\right].
\end{align}

Therefore

\begin{align}
\kappa\to0:
\qquad
a_{1,a}=a_{2,a}=0,
\end{align}

while bounded energy in the strong-gluing limit gives

\begin{align}
\kappa\to\infty:
\qquad
a_{1,a}=a_{2,a}.
\end{align}

The common dressed field becomes unconstrained by the vanishing $\kappa^{-1}$ coefficient, while the relative dressed field is suppressed by the growing $2\kappa$ coefficient. The finite-$\kappa$ cut equations then match the normal field strengths. After pulling back the CPS form and quotienting its null directions, the bounded-data limit gives the Maxwell phase space on the uncut spacetime.

The last statement assumes that the unscaled gauge-invariant bulk field strengths and dressed cut fields remain uniformly bounded as $\kappa\to\infty$. It is not the only possible scaling limit. A CPS-normalized family can have a common boundary component of order $\kappa^{1/2}$; after rescaling, that component survives as a massless interface scalar. Thus the unscaled bounded-data limit gives pure uncut Maxwell, whereas the normalized-mode scaling gives uncut Maxwell together with an additional interface theory. These are different limiting phase spaces, not two descriptions of the same limit.

At finite mode cutoff, the two boundary eigenchannels again require independent coefficients. After dualization their oscillator pencils have the form

\begin{align}
K_N=D_N,
\qquad
G_{s,N}=I+\frac1{\alpha_s}bb^{\mathrm T},
\qquad
\alpha_s\in\{d,d+2h\}.
\end{align}

If $T_N(k)$ is the omitted boundary response, exact matching at a selected wave number requires

\begin{align}
\alpha_{s,N}(k)=\alpha_s-k^2T_N(k).
\end{align}

The correction depends on both the boundary channel and the selected energy. Consequently, even allowing two independent constants $d_N$ and $h_N$ does not exactly match more than one generic spectral point in each channel. A finite energy window requires energy-dependent boundary operators or an equivalent auxiliary-mode completion.

\subsubsection{Compact gauge group}\label{compact-gauge-group}

For compact $U(1)$, each $\varphi_i$ may be circle-valued,

\begin{align}
\varphi_i(t,y+\ell_y)
=\varphi_i(t,y)+2\pi s_i,
\qquad
s_i\in\mathbb Z.
\end{align}

If the physical outer holonomy is fixed to zero, strict Dirichlet data imply

\begin{align}
\oint_{S_y^1}A_{i,y}\big|_{\Gamma_i}\,\mathrm dy=2\pi s_i,
\qquad
\int_{\Sigma_i}\mathrm dx\mathrm dy\,B_i=2\pi n_i s_i.
\end{align}

The integers label disconnected magnetic-flux components of the strict Dirichlet region theories. Under strong gluing, the relative winding $s_1-s_2$ becomes the magnetic-flux label of the global connection, whereas simultaneous changes of the two cut trivializations do not create an additional global sector. These discrete components are separate from the positive-frequency photon modes and must be included when specifying the compact theory's full phase space.
